\section{Theoretical Justification for Sensitivity-Constrained Training}
\label{sec_sup_method}

\renewcommand{\thefigure}{A.\arabic{figure}}  % Number figures as A.1, A.2, ...
\renewcommand{\thetable}{A.\arabic{table}}  % Number tables as A.1, A.2, ...


Neural Operators (NOs) provide a framework for learning mappings between infinite-dimensional function spaces, enabling resolution-invariant PDE approximations \citep{lu2021learning, kovachki2023neural}. Unlike traditional solvers that discretize PDEs explicitly, NOs approximate a solution operator \( \mathcal{G}: \mathcal{U} \to \mathcal{V} \), where \( \mathcal{U} \) represents the space of input functions (e.g., initial conditions, boundary conditions, parameters), and \( \mathcal{V} \) denotes the space of solutions. Given an input function \( a(x, t, p) \in \mathcal{U} \), which encapsulates problem-specific information, and an output solution \( u(x, t, p) \in \mathcal{V} \), the operator mapping is:

\begin{equation}
u(x, t, p) = \mathcal{G}(a)(x, t, p),
\end{equation}

where \( x \in \mathcal{X} \) are spatial coordinates, \( t \in [0, T] \) represents time, and \( p \in \mathbb{R}^m \) denotes system parameters. Neural Operators parameterize \( \mathcal{G} \) with learnable weights \( \theta \), forming \( \mathcal{G}_\theta \):

\begin{equation}
u(x, t, p) = \mathcal{G}_\theta(a)(x, t, p).
\end{equation}

Unlike conventional neural networks, NOs operate on function spaces, allowing solutions to be transferred across different resolutions.

\subsection*{Neural Operators}
The architecture of a Neural Operator consists of iterative layers that update the latent representation \( v_t(x, t, p) \) via:

\begin{equation}
v_{t+1}(x, t, p) = \sigma \left( \mathcal{K}_\theta(v_t)(x, t, p) + \mathcal{L}_\theta(v_t)(x, t, p) \right),
\end{equation}

where \( \mathcal{K}_\theta \) is a nonlocal integral operator, \( \mathcal{L}_\theta \) is a local transformation, and \( \sigma \) is a nonlinear activation function. The nonlocal operator is defined as:

\begin{equation}
\mathcal{K}_\theta(v_t)(x, t, p) = \int_{\mathcal{X}} \int_0^T \int_{\mathbb{R}^m} K_\theta(x, t, p; y, s, q) v_t(y, s, q) \, dy \, ds \, dq,
\end{equation}

where \( K_\theta \) is a learnable kernel encoding interactions across space, time, and parameters. The local transformation follows:

\begin{equation}
\mathcal{L}_\theta(v_t)(x, t, p) = W_\theta v_t(x, t, p) + b_\theta(x, t, p).
\end{equation}

Different parameterization strategies yield distinct NO architectures:

- \textbf{Fourier Neural Operators (FNOs)} employ spectral transforms for efficiency, defining the integral operator in Fourier space:

\begin{equation}
\mathcal{F}(\mathcal{K}_\theta(v_t))(k, \omega) = R_\theta(k, \omega) \cdot \mathcal{F}(v_t)(k, \omega),
\end{equation}

where \( \mathcal{F} \) denotes the Fourier transform, and \( R_\theta(k, \omega) \) is a learnable kernel truncated for high frequencies:

\begin{equation}
R_\theta(k, \omega) = 0 \quad \text{for} \quad |k|, |\omega| > k_{\text{max}}.
\end{equation}

- \textbf{Wavelet Neural Operators (WNOs)} replace Fourier transforms with wavelet decompositions to capture localized structures across scales \citep{tripura2023wavelet}.

- \textbf{DeepONets} decompose the mapping into a trunk network (spatial-temporal encoding) and a branch network (parametric encoding) \citep{lu2021learning}:

\begin{equation}
\mathcal{K}_\theta(v_t)(x, t, p) = \sum_{i=1}^m b_i(p) \tau_i(x, t),
\end{equation}

where \( b_i(p) \) and \( \tau_i(x, t) \) are learned functions.

% - Convolutional Neural Operators (CNOs) use convolutional layers for localized information transfer, making them efficient for structured grid-based PDEs \citep{raonic2023convolutional}.

After \( T \) iterations, the final output is projected onto the solution space:

\begin{equation}
u(x, t, p) = Q_\theta(v_T)(x, t, p).
\end{equation}

\subsection*{Sensitivity-Constrained Neural Operators (SC-NO)}
Standard NOs minimize a solution-based loss:

\begin{equation}
L_u = \ell_u(\hat{u}(x, t; p), u(x, t; p)),
\end{equation}

where \( \ell_u \) measures prediction error. However, they do not explicitly account for how solutions change with respect to parameters \( p \), which is crucial for generalization and inverse problems. To address this, we introduce a sensitivity regularization term:

\begin{equation}
L_S = \ell_s\left( \frac{\partial \hat{u}(x, t; p)}{\partial p}, \frac{\partial u(x, t; p)}{\partial p} \right),
\end{equation}

where \( \frac{\partial \hat{u}}{\partial p} \) is computed via automatic differentiation, and \( \frac{\partial u}{\partial p} \) is derived from a high-fidelity solver. The final training objective balances both terms:

\begin{equation}
L = L_u + \lambda L_S,
\end{equation}

where \( \lambda \) controls sensitivity enforcement. 






\subsection*{Theoretical Justification} \label{prove}

We provide a mathematical basis for why Sensitivity-Constrained Neural Operators (SC-NOs) outperform standard training by incorporating sensitivity loss, improving generalization, long-term stability, and inversion accuracy. We define the operator $\mathcal{G}_\theta: \mathcal{U} \to \mathcal{V}$, which maps an input function $a(x) = (u_0(x), p(x))$ to $\hat{u}(x,t) = \mathcal{G}_\theta(a(x))$, approximating the true solution $u(x,t)$.

\subsubsection*{Assumptions and Notation}

We work in bounded domains: the parameter space $\mathcal{P}_{\text{param}} \subset \mathbb{R}^m$ and the spatio-temporal domain $\mathcal{X} \times [0,T]$. We assume $\hat{u}$ is differentiable with respect to $p$. We use the Euclidean norm $|\cdot|$ on $\mathbb{R}^m$ and the corresponding induced operator norms. We assume the true solution $u$ and its sensitivity $\frac{\partial u}{\partial p}$ are bounded by $M$ and $M_p$ respectively, for all $(x,t,p)$.
\begin{equation}
|u(x,t;p)| \leq M, \quad \left\| \frac{\partial u}{\partial p}(x,t;p) \right\| \leq M_p.
\end{equation}

Our core argument relies on standard (though strong) assumptions from learning theory: for an overparameterized, well-trained model $\theta_s^*$, a small empirical loss implies a small uniform error bound.

\paragraph{Assumption 1 (Uniform Convergence of Gradients):}
We assume that a model $\theta_s^*$ trained to $L_s(\theta_s^*)$ on $n$ samples satisfies:
\begin{equation}
\sup_{x,t,p} \left\| \frac{\partial \hat{u}}{\partial p}(x,t;p;\theta_s^*) - \frac{\partial u}{\partial p}(x,t;p) \right\| \le \varepsilon_s
\end{equation}
where $\varepsilon_s \to 0$ as $L_s \to 0$ and $n \to \infty$.

\paragraph{Assumption 2 (Uniform Convergence of Solutions):}
We assume a similar uniform bound $\varepsilon_u$ for the solution itself:
\begin{equation}
\sup_{x,t,p} |\hat{u}(x,t;p;\theta_s^*) - u(x,t;p)| \le \varepsilon_u
\end{equation}
where $\varepsilon_u \to 0$ as $L_u \to 0$ and $n \to \infty$.

\subsubsection*{Cost Functionals}

Training minimizes the following losses, where $\frac{\partial u}{\partial p}$ is computed on-the-fly via a differentiable solver or adjoint model:
\begin{equation}
L_u(\theta) = \frac{1}{n} \sum_{i=1}^n [\hat{u}(x_i, t_i; p_i;\theta) - u(x_i, t_i; p_i)]^2,
\end{equation}
\begin{equation}
L_s(\theta) = \frac{1}{n} \sum_{i=1}^n \left\| \frac{\partial \hat{u}}{\partial p}(x_i, t_i; p_i;\theta) - \frac{\partial u}{\partial p}(x_i, t_i; p_i) \right\|^2,
\end{equation}
\begin{equation}
L(\theta) = L_u(\theta) + \lambda L_s(\theta), \quad \lambda > 0.
\end{equation}


\subsubsection*{Improved Generalization}

Training with $L(\theta)$ acts as a powerful regularizer by explicitly constraining the model's Lipschitz constant with respect to parameters. We define this constant as:
\begin{equation}
L_{\text{lip}}(\theta) = \sup_{x,t,p} \left\| \frac{\partial \hat{u}}{\partial p}(x,t;p;\theta) \right\|.
\end{equation}
For the sensitivity-constrained model $\theta_s^*$, we can rigorously bound its Lipschitz constant using Assumption 1 and the triangle inequality:
\begin{align}
L_{\text{lip}}(\theta_s^*) &= \sup \left\| \frac{\partial u}{\partial p} + \left( \frac{\partial \hat{u}}{\partial p} - \frac{\partial u}{\partial p} \right) \right\| \notag \\
&\le \sup \left\| \frac{\partial u}{\partial p} \right\| + \sup \left\| \frac{\partial \hat{u}}{\partial p} - \frac{\partial u}{\partial p} \right\| \notag \\
&\le M_p + \varepsilon_s.
\end{align}
By driving $\varepsilon_s \to 0$, we force the model to inherit the (bounded) Lipschitz constant of the true physical system. This has a direct impact on generalization. For fixed $(x,t)$, we define the hypothesis class $\mathcal{F}_{\text{SC}} = \{p \mapsto \hat{u}(x,t;p;\theta): L_{\text{lip}}(\theta) \le M_p+\varepsilon_s\}$. On a bounded parameter domain, standard results imply that the empirical Rademacher complexity $\hat{R}_n(\mathcal{F}_{\text{SC}})$ of this class scales linearly with its Lipschitz constant, $\hat{R}_n(\mathcal{F}_{\text{SC}}) \lesssim (M_p + \varepsilon_s)$. This complexity is strictly lower than that of an unconstrained class $\mathcal{F}_u$ that admits arbitrarily large Lipschitz constants. The standard generalization bound,
\begin{equation}
\mathbb{E}[(\hat{u} - u)^2] \leq L_u(\theta) + 2 \hat{R}_n(\mathcal{F}) + \text{ComplexityPenalty}(\delta, n),
\end{equation}
is therefore tighter for $\mathcal{F}_{\text{SC}}$, implying better performance on unseen parameters.


\subsubsection*{Long-Term Prediction Stability}

This argument is most clearly made by setting the parameter $p$ to be the initial condition $u_0$. The sensitivity $\frac{\partial u(t)}{\partial u_0}$ is the \textbf{fundamental solution operator}, and its norm governs the growth of perturbations (i.e., the Lyapunov exponents).

Assume the true physical system's growth is bounded:
\begin{equation}
\left\| \frac{\partial u(t)}{\partial u_0} \right\| \le C e^{\alpha t}
\end{equation}
for some physical exponent $\alpha$ (which could be positive for a chaotic system, or negative for a stable one). A standard model $\theta_u^*$, trained only on $L_u$, may learn an incorrect operator $\frac{\partial \hat{u}}{\partial u_0}$ with a much larger exponent $\hat{\alpha} \gg \alpha$, leading to spurious amplification of small errors. The SC-NO, by minimizing $L_s$ (with $p=u_0$), uses Assumption 1 to force its learned operator to be close to the true one:
\begin{align}
\left\| \frac{\partial \hat{u}(t)}{\partial u_0} \right\| &\le \left\| \frac{\partial u(t)}{\partial u_0} \right\| + \sup \left\| \frac{\partial \hat{u}}{\partial u_0} - \frac{\partial u}{\partial u_0} \right\| \notag \\
&\le C e^{\alpha t} + \varepsilon_s.
\end{align}
This shows that the SC-NO reproduces the physical growth rates of perturbations up to $O(\varepsilon_s)$ and does not introduce spurious instabilities beyond those present in the governing system.


\subsubsection*{Inversion Accuracy}

In parameter inversion, we seek to find $p^*$ by minimizing a loss $\ell$ between our model $\hat{u}$ and an observation $u_{\text{obs}}$. Let $J(p) = \ell(u(p), u_{\text{obs}})$ be the "true" (but intractable) loss landscape, and $J_\theta(p) = \ell(\hat{u}_\theta(p), u_{\text{obs}})$ be the surrogate landscape our optimizer actually sees.

Gradient descent updates use the gradient of $J_\theta$:
\begin{equation}
p_{k+1} = p_k - \eta \nabla_p J_\theta(p_k), \quad \text{where} \quad \nabla_p J_\theta = \frac{\partial \ell}{\partial \hat{u}} \cdot \frac{\partial \hat{u}}{\partial p}.
\end{equation}
The success of the inversion depends on the \textbf{gradient error} $\|\nabla J_\theta(p) - \nabla J(p)\|$. If $\ell$ is smooth, this error is bounded by our assumed uniform errors:
\begin{equation}
\|\nabla J_\theta(p) - \nabla J(p)\| \le C_1 \varepsilon_u + C_2 \varepsilon_s
\end{equation}
for some constants $C_1, C_2$.

A standard model $\theta_u^*$ is only trained to minimize $\varepsilon_u$, leaving a large, uncontrolled gradient error $C_2 \varepsilon_s$. The optimizer is fed an inaccurate gradient. An SC-NO is explicitly trained to drive both $\varepsilon_u \to 0$ and $\varepsilon_s \to 0$. We further assume that, in a neighborhood of $p^*$, the true objective $J(p)$ is $\mu$-strongly convex. Under this condition, standard perturbation results for gradient descent imply that optimization on $J_\theta$ converges to an $O(\varepsilon_u + \varepsilon_s)$ neighborhood of the true optimum $p^*$. By minimizing $\varepsilon_s$, SC-NOs guarantee a high-fidelity gradient, ensuring a much more accurate and stable convergence.







%\subsection*{Jacobian Computation}
%The accuracy of Sensitivity-Constrained Neural Operators (SC-NOs) relies on aligning the predicted solution \(\hat{u}(x, t, p) = \mathcal{G}_\theta(a)(x, t, p)\) and its sensitivity \(\frac{\partial \hat{u}}{\partial p}\) with the true solution \( u(x, t, p) \) and sensitivity \(\frac{\partial u}{\partial p}\). During data preparation, we obtain true Jacobians using two strategies: (1) differentiable solvers with gradient tracing, where automatic differentiation is enabled during simulation, and (2) adjoint methods, which efficiently compute gradients using a \textit{Discretize-then-Optimize} approach. Both methods allow us to extract high-fidelity sensitivity information while maintaining computational efficiency. During training, computing the full Jacobian of \(\hat{u}\) at every spatial point remains intractable. To address this, we employ gradient sub-sampling, where automatic differentiation is used to compute \(\frac{\partial \hat{u}}{\partial p}\) at a randomly selected subset of \( n \) spatial locations (\( n < n_x \times n_y \)) in each iteration. These sampled gradients are compared against precomputed true gradients obtained from the solvers. Since different sampling points are selected in each training step, the entire spatial domain is gradually covered over multiple iterations, ensuring comprehensive sensitivity alignment while avoiding excessive memory consumption. This approach allows SC-NOs to enforce sensitivity constraints effectively without the prohibitive cost of full Jacobian computation, making them scalable for high-dimensional and computationally demanding problems. While our implementation is on PyTorch, it can be similarly done on other platforms like JAX, Julia and Tensorflow, or conventional modeling environment like Fortran or C where adjoints are implemented.

